% !PRINTED 297
\weberpage{297}
By the concept of the normal divisor, $\mathfrak{N}$ also contains the substitution
\weq{12}{\begin{pmatrix} 0, & -1 \\ 1, & 0 \end{pmatrix}\begin{pmatrix} 0, & 1 \\ -1, & \xi \end{pmatrix}\begin{pmatrix} 0, & 1 \\ -1, & 0 \end{pmatrix} = \begin{pmatrix} \xi, & 1 \\ -1, & 0 \end{pmatrix}}
and consequently also
\weq{13}{\begin{pmatrix} 0, & 1 \\ -1, & \xi \end{pmatrix}\begin{pmatrix} \xi, & 1 \\ -1, & 0 \end{pmatrix} = \begin{pmatrix} 1, & 0 \\ 2\xi, & 1 \end{pmatrix} \equiv A^{2\xi}\ [\S\ 79,\ (11)],}
hence also $A$ and all its powers. Therefore $\mathfrak{N}$ also contains
\weq{14}{\begin{pmatrix} 1, & 0 \\ \eta, & 1 \end{pmatrix}\begin{pmatrix} 0, & 1 \\ -1, & \xi \end{pmatrix} = \begin{pmatrix} 0, & 1 \\ -1, & \xi+\eta \end{pmatrix}}
for an arbitrary $\eta$, that is, the substitution (11) for \gesperrt{every arbitrary} $\xi$, and consequently also
\[
B = \begin{pmatrix} 0, & 1 \\ -1, & 0 \end{pmatrix}.
\]
Accordingly, by the theorem of the preceding section, the group $\mathfrak{N}$ is identical with $\mathfrak{L}_0$.

It still remains to discuss the possibility that in (11) $\xi \equiv 0$.

In this case, then, the group $\mathfrak{N}$ contains the substitution
\[
B = \begin{pmatrix} 0, & 1 \\ -1, & 0 \end{pmatrix},
\]
and consequently also, for an arbitrary $\alpha$ not divisible by $p$,
\[
V = \begin{pmatrix} 0, & \alpha \\ -\alpha^{-1}, & 0 \end{pmatrix}\begin{pmatrix} 0, & 1 \\ -1, & 0 \end{pmatrix}\begin{pmatrix} 0, & -\alpha \\ \alpha^{-1}, & 0 \end{pmatrix}\begin{pmatrix} 0, & -1 \\ 1, & 0 \end{pmatrix} = \begin{pmatrix} \alpha^2, & 0 \\ 0, & \alpha^{-2} \end{pmatrix}.
\]
From this one derives further, as contained in $\mathfrak{N}$:
\[
W = \begin{pmatrix} 1, & 0 \\ -1, & 1 \end{pmatrix}\begin{pmatrix} \alpha^{-2}, & 0 \\ 0, & \alpha^2 \end{pmatrix}\begin{pmatrix} 1, & 0 \\ 1, & 1 \end{pmatrix}\begin{pmatrix} \alpha^2, & 0 \\ 0, & \alpha^{-2} \end{pmatrix} = \begin{pmatrix} 1, & 0 \\ \alpha^4-1, & 1 \end{pmatrix}.
\]
If now $p$ is greater than 5, then $\alpha$ can be so assumed that $\alpha^4 - 1$ is not divisible by $p$; then $\mathfrak{N}$ also contains the substitution $A$ and its powers, and consequently $\mathfrak{N}$ is identical with $\mathfrak{L}_0$.

If $p = 5$, then $\alpha^4 - 1$ is always divisible by $p$, so this conclusion is not applicable.

In this case, however, it follows from the concept of the normal divisor, as contained in $\mathfrak{N}$:
%CONT>>
